%Schriftgroesse, Layout, Papierformat, Art des Dokumentes
\documentclass[listof=totoc, 12pt, oneside, bibliography=totoc, a4paper, parskip=full+, final, headinclude=false, footinclude=false]{scrreprt} 

%Einstellungen der Seitenraender, Zeilenabstand
\usepackage[left=4cm,right=3cm,top=3cm,bottom=2cm,bindingoffset=0cm]{geometry}
\usepackage{titlesec}                  % Ueberschriftenformatierungen 
\usepackage{setspace} 
\setlength{\footskip}{0.8cm}
%bestimmt abstand zwischen footnotes und text
\setlength{\skip\footins}{6mm}

\usepackage{pdfpages}
%anpassung des Abstadnes zwischen zwei Footnotes
\setlength {\footnotesep}{0mm} 
%f�r subusubsection Nummerierungen
\setcounter{tocdepth}{3}
\setcounter{secnumdepth}{3} 

\usepackage{trfsigns}
%F�r timing diagramme
\usepackage{tikz-timing}[2009/05/15]
\def\degr{${}^\circ$}

%% Zeilenabstand
\onehalfspacing                                     % Zeilenabstand 1,5 fach

%Deutsche Rechtschreibung f�r Text und Lit.-verzeichnis
\usepackage[ngerman]{babel}
\usepackage[german]{babelbib} 
%zur korrekten Schreiben vn Anf�hrungszeichen
\usepackage[babel,german=quotes]{csquotes}
%Schriftart �ndern 
\usepackage{mathptmx}
%pacakge f�r mathematische Transormationszeichen
\usepackage{trfsigns}
%BiblatexStyle alphadin = name+Erscheinungsjahr
\bibliographystyle{unsrt}

%bibs als Section
\makeatletter
\renewcommand*\bib@heading{%
	\section*{\bibname}%
	\@mkboth{\leftmark}{\rightmark}%
}
\makeatother
\usepackage{relsize}
\usepackage{lipsum}

% Ver�ndert Bildunterschriften
\addtokomafont{caption}{\small \it} 

%Kein Einr�cken bei neuem Absatz
\setlength\parindent{0pt} 

%f�r Abk�rzungsverzeichnis
\usepackage{acronym}

%Fuer die Bibliothek
\usepackage{cite}

\usepackage[center]{caption}

% F�r Symbole
\usepackage{textcomp}

% Damit einzelne Seiten im Querformat angezeigt werden k�nnen
\usepackage{pdflscape}

%Text in meheren Spalten
\usepackage{multicol}
%Um Tabellen in der tabular Umgebung zu nutzen
\usepackage{tabularx}

%Mathematische Zeichen wie Menge, Element usw
\usepackage{dsfont}

%for Tabulares
\newcolumntype{C}[1]{>{\centering\arraybackslash}p{#1}}

%Umlaute ermoeglichen
\usepackage[utf8]{inputenc}
\usepackage[T1]{fontenc}
\usepackage{lmodern}
\usepackage{textcomp}

%Fuer Grafiken
\usepackage{graphicx}
\usepackage{svg}

%Um mehrere Bilder nebeneinander darstellen zu koennen
\usepackage{subfigure}

%f�r umrandete Texte -> Sperrvermerk
\usepackage{framed}

%Um Grafiken neben dem Text einzufuegen
\usepackage{wrapfig}

%Fuer die eingabe mathematischer Formeln
\usepackage{amsmath, amsthm, amssymb}
\usepackage{mathtools}

%Bildunterschriften manipulieren
\captionsetup{format = plain}
\addto\captionsngerman
{
	\renewcommand{\figurename}{Abb.} %Fuer die Bildunterschriften wird nun Abb. x verwendet
	\renewcommand{\tablename}{Tab.} 
}

%Kopf- und Fusszeile bzw eigene Designs, hier sind �nderungen einzutragen
\usepackage[automark]{scrlayer-scrpage}

\newpagestyle{WissDokuNorm}
{ %
	(0pt,0pt)
	{\hfill\headmark}
	{\headmark\hfill}
	{\rlap{}\hfill%
		\llap{\headmark\hfill}}
	(\textwidth,.4pt)
}
{
	(\textwidth,.4pt)
	{\pagemark\hfill}
	{\hfill\pagemark}
	{Mobile Computing Final Report - LogChirpy\hfill\pagemark}
	(0pt,0pt)
}

\newpagestyle{WissDokuChapter}
{ %
	(0pt,0pt)
	{\hfill}
	{\hfill}
	{\rlap{}\hfill%
		\llap{\hfill}}
	(0pt,.4pt)
}{%
(\textwidth,.4pt)
{\pagemark\hfill}
{\hfill\pagemark}
{Mobile Computing Final Report - LogChirpy\hfill\pagemark}
(0pt,0pt)
}

\newpagestyle{WissDokuNothing}
{ %
	(0pt,0pt)
	{\hfill}
	{\hfill}
	{\rlap{}\hfill%
		\llap{\hfill}}
	(0pt,.4pt)
}{%
(0pt,0pt)
{\pagemark\hfill}
{\hfill\pagemark}
{\hfill\pagemark}
(0pt,0pt)
}

%Zuweisungen der gew�nschten styles
\assignpagestyle{\chapter}{WissDokuChapter}
\assignpagestyle{\section}{WissDokuNorm}
\assignpagestyle{empty}{WissDokuNothing}

\pagestyle{WissDokuNorm}

%----------------------------------------------

%F�r korrektes anzeigen des Quellcodes
\usepackage[utf8]{inputenc}
\usepackage{color}
\definecolor{darkblue}{rgb}{0,0,.6}
\definecolor{darkred}{rgb}{.6,0,0}
\definecolor{darkgreen}{rgb}{0,.6,0}
\definecolor{lightblue}{rgb}{0.97,0.99,1}

\usepackage{listings}
\renewcommand{\lstlistingname}{Codeauszug}
\lstdefinestyle{javascript}{language=JavaScript, frame=lr, morekeywords={const,let,var,function,class,async,await,export,import,from,React,useState,useEffect}}
\lstdefinestyle{typescript}{language=Java, frame=lr, morekeywords={interface,type,const,let,var,function,class,async,await,export,import,from,React,useState,useEffect}}
\lstset
{
	style=javascript,
	basicstyle=\ttfamily\small,
	commentstyle=\color{darkgreen},
	keywordstyle=\bfseries\color{darkblue},
	stringstyle=\color{darkred},
	showspaces=false,
	showstringspaces=false,
	showtabs=false,
	columns=fixed,
	frame=single,
	frameround=ffff,
	numbers=left,
	numberstyle=\tiny,
	numbersep=5pt,
	breaklines=true,
	backgroundcolor=\color{lightblue},
	captionpos=b
}

%l�dt das Paket zum erzwingen der Grafikposition
\usepackage{here} 

% erm�glich es, mit \ScaleIfNeeded Grafik nur dann skalieren zu lassen, wenn sie gr��er als eine seite sind
\makeatletter
\def\ScaleIfNeeded{%
	\ifdim\Gin@nat@width>\linewidth
	\linewidth
	\else
	\Gin@nat@width
	\fi
}
\makeatother

% Keine "Schusterjungen"
\clubpenalty = 10000

% Keine "Hurenkinder"
\widowpenalty = 10000
\displaywidowpenalty = 10000

\raggedbottom

\setlength{\parskip}{3ex plus 3ex minus 3ex}

\renewcommand*\chapterheadstartvskip{\vspace{-\topskip}}
\renewcommand*{\chapterheadendvskip}{\vspace*{0pt}}      % Abstand zwischen Chapter u. Fliesstext

%Um Links im Dokument zu aktivieren
\usepackage{hyperref}
\hypersetup{
	colorlinks=false,
	pdfborder={0 0 0},
}

%Autor und Titel
\author{Martin Lauterbach, Luis Wehrberger, Youmna Samouneh}
\title{LogChirpy - Ornithological Archival App}